% ============================================================
% CE 749 — Project Report Template
% Bi-Objective MCF for Indian Railway Freight Allocation
% Indian Institute of Technology Bombay
% ============================================================

\documentclass[11pt, a4paper]{article}

% ── Packages ─────────────────────────────────────────────────
\usepackage[a4paper, top=0.8in, bottom=0.8in, left=1in, right=1in]{geometry}
\usepackage{times}                      % Times New Roman font
\usepackage{setspace}                   % Line spacing
\usepackage{titlesec}                   % Section heading formatting
\usepackage{fancyhdr}                   % Header and footer
\usepackage{graphicx}                   % Figures
\usepackage{float}                      % Figure placement [H]
\usepackage{caption}                    % Caption formatting
\usepackage{subcaption}                 % Subfigures
\usepackage{booktabs}                   % Professional tables
\usepackage{array}                      % Table column formatting
\usepackage{longtable}                  % Multi-page tables
\usepackage{multirow}                   % Merged table rows
\usepackage{amsmath}                    % Math equations
\usepackage{amssymb}                    % Math symbols
\usepackage{mathtools}                  % Extended math
\usepackage{bm}                         % Bold math
\usepackage{algorithm}                  % Algorithm block
\usepackage{algpseudocode}              % Pseudocode
\usepackage[hidelinks]{hyperref}        % Clickable links (no red boxes)
\usepackage{url}                        % URL formatting
\usepackage[numbers, sort&compress]{natbib}   % Bibliography
\usepackage{enumitem}                   % List formatting
\usepackage{xcolor}                     % Colours
\usepackage{tikz}                       % Diagrams
\usepackage{pgfplots}                   % Plots in LaTeX
\pgfplotsset{compat=1.18}
\usepackage{indentfirst}                % Indent first paragraph
\usepackage{parskip}                    % Paragraph spacing

% ── Colours ──────────────────────────────────────────────────
\definecolor{iitbblue}{RGB}{0, 75, 135}
\definecolor{darkgray}{RGB}{80, 80, 80}
\definecolor{lightgray}{RGB}{240, 240, 240}

% ── Page Layout ──────────────────────────────────────────────
\setstretch{1.5}                        % 1.5 line spacing
\setlength{\parindent}{0.5cm}
\setlength{\parskip}{6pt}

% ── Section Heading Styles ───────────────────────────────────
\titleformat{\section}
  {\normalfont\fontsize{14}{17}\bfseries\color{iitbblue}}
  {\thesection}{1em}{}
  [\titlerule]

\titleformat{\subsection}
  {\normalfont\fontsize{12}{15}\bfseries}
  {\thesubsection}{1em}{}

\titleformat{\subsubsection}
  {\normalfont\fontsize{12}{15}\bfseries\itshape}
  {\thesubsubsection}{1em}{}

\titlespacing*{\section}{0pt}{18pt}{8pt}
\titlespacing*{\subsection}{0pt}{12pt}{4pt}
\titlespacing*{\subsubsection}{0pt}{10pt}{4pt}

% ── Header / Footer ──────────────────────────────────────────
\pagestyle{fancy}
\fancyhf{}
\renewcommand{\headrulewidth}{0.4pt}
\renewcommand{\footrulewidth}{0.4pt}
\fancyhead[L]{\small\textit{CE 749 : Freight Transportation Planning and Logistics}}
\fancyhead[R]{\small\textit{IIT Bombay}}
\fancyfoot[C]{\small\thepage}

% ── Caption Formatting ───────────────────────────────────────
\captionsetup{
  font=small,
  labelfont=bf,
  labelsep=period,
  justification=centering
}

% ── Equation Numbering by Section ────────────────────────────
\numberwithin{equation}{section}
\numberwithin{figure}{section}
\numberwithin{table}{section}

% ── Hyperlink Setup ──────────────────────────────────────────
\hypersetup{
    pdftitle={CE749 Project Report},
    pdfauthor={Niraj — 25M1528},
    colorlinks=false
}

% ============================================================
%  DOCUMENT BEGINS
% ============================================================
\begin{document}

% ── Title Page ───────────────────────────────────────────────
\begin{titlepage}
    \begin{center}
        \vspace*{0.5cm}

        {\Large \textbf{Indian Institute of Technology Bombay}} \\[0.3cm]
        
        \rule{\linewidth}{0.5pt} \\[0.5cm]

        {\large\textit{CE 749 : Freight Transportation Planning and Logistics}} \\[0.3cm]
        {\large \textit{Course Project Report}} \\[1.5cm]

        {\LARGE \textbf{Bi-Objective Optimization of Freight Allocation}} \\[0.4cm]
        {\LARGE \textbf{on the Indian Railway Network}} \\[0.8cm]
        {\Large \textbf{(A Multi-Commodity Flow Approach with}} \\[0.3cm]
        {\Large \textbf{Transshipment, Backhaul, and Green Routing)}} \\[2.0cm]

        \rule{\linewidth}{0.5pt} \\[0.8cm]

        \begin{tabular}{rl}
            \textbf{Submitted by:}  & Niraj  \\
            \textbf{Roll No:}       & 25M1528 \\
            \textbf{Submitted to:}  & Prof. Gopal Patil \\
            \textbf{Semester:}          & Spring 2026 \\
        \end{tabular} \\[2.0cm]

        \rule{\linewidth}{0.5pt}

        \vfill
        {\small Indian Institute of Technology Bombay, Powai, Mumbai -- 400076}
    \end{center}
\end{titlepage}

% ── Abstract ─────────────────────────────────────────────────
\newpage
\pagenumbering{roman}
\setcounter{page}{1}

\begin{abstract}
\addcontentsline{toc}{section}{Abstract}
\noindent
Indian Railways transported 1,509~MT of freight in FY~2022-23, with the Golden
Quadrilateral corridor carrying 55\% of this volume over just 16\% of total route
length. The recent commissioning of the Eastern and Western Dedicated Freight Corridors
has introduced fully electrified, high-capacity alternatives alongside existing diesel
conventional lines, creating a structural conflict between cost minimisation and
CO\textsubscript{2} emission reduction in freight routing decisions.

This report formulates a \textbf{Bi-Objective Multi-Commodity Flow Linear Program
(MCF-LP)} for annual freight allocation on the Indian Railways Golden Quadrilateral
network. The model routes four commodity types — Coal, Cement, Foodgrains, and
Fertilizers — across a 16-node, 50-arc network representing the GQ and its diagonals,
with a total annual demand of 435~MT across 28 origin-destination pairs. Three
operationally realistic features distinguish the formulation from standard MCF models:
(i)~transshipment costs at all 12 intermediate junction nodes capturing marshalling yard
detention charges; (ii)~empty wagon backhaul repositioning costs modelled as a linear
auxiliary objective component; and (iii)~traction-differentiated emission factors
(electric: 0.010~kgCO\textsubscript{2}/tonne-km; diesel: 0.035~kgCO\textsubscript{2}/tonne-km)
sourced from the India Default GHG Emission Values report (SFC India, 2025). The two
objectives — total logistics cost ($Z_1$) and total CO\textsubscript{2} emissions
($Z_2$) — are optimised simultaneously using the \textbf{epsilon-constraint method},
which generates a 12-point Pareto frontier.

Results show that the emission-optimal solution achieves an \textbf{8.66\% reduction
in CO\textsubscript{2}} (from 10.22~Mt to 9.34~Mt) at a \textbf{cost premium of
15.45\%} over the cost-optimal baseline of Rs~99,614~crore. Empty wagon backhaul repositioning accounts for \textbf{25.8\% of total
logistics cost} — a component absent from standard single-term MCF formulations.
A Greedy Shortest Path heuristic benchmark shows that naive routing produces solutions
that cost \textbf{10.22\% more} and emit \textbf{9.69\% more} CO\textsubscript{2} than the MCF-LP optimum, confirming
the value of formal optimisation for both cost and environmental performance. Sensitivity analysis
identifies diesel arc capacity of 70~MT as the binding constraint threshold, beyond
which cost plateaus completely. The carbon price breakeven is estimated at
Rs~1,73,930/tonne~CO\textsubscript{2}, providing a concrete threshold for green routing
policy on Indian Railways.

\end{abstract}

\vspace{1cm}

\noindent\textbf{Keywords:} Multi-commodity flow, bi-objective optimization, Indian Railways, 
epsilon-constraint, green freight, transshipment, backhaul, Pareto frontier, 
Golden Quadrilateral, CO\textsubscript{2} emissions.

% ── Table of Contents ────────────────────────────────────────
\newpage
\tableofcontents

\newpage
\listoffigures
\listoftables

% ── Main Body ────────────────────────────────────────────────
\newpage
\pagenumbering{arabic}
\setcounter{page}{1}

% ============================================================
\section{Introduction}
\label{sec:intro}
% ============================================================

Indian Railways is one of the largest freight rail networks in the world, transporting
1,509.10 million tonnes (MT) of goods in FY~2022-23 and generating freight revenue of
Rs~1,60,158~crore \cite{iryearbook2023}. With a route network spanning over 67,000~km and
a fleet of more than 3~lakh wagons, it serves as the circulatory system of India's
industrial economy, moving coal to power plants, cement to construction sites,
foodgrains through the Food Corporation of India supply chain, and fertilizers to
agricultural distribution hubs across the country.

Despite this scale, a critical planning challenge persists at the operational level: given
a set of freight demands across multiple origin-destination (OD) pairs and multiple
commodity types, \textit{how should freight be allocated across the railway network so as
to minimise total logistics cost}? This question is the freight allocation problem — and it is not easy to solve. Every year, suboptimal routing decisions translate directly into
unnecessary expenditure on longer hauls, avoidable marshalling yard detention charges,
and inefficient empty wagon repositioning. Conversely, the growing policy emphasis on
decarbonising freight transport introduces a second, competing objective: minimising
CO\textsubscript{2} emissions. These two objectives are in structural conflict on the
Indian Railways network, where lower-cost routes frequently pass through diesel-traction
corridors that emit up to 3.5 times more CO\textsubscript{2} per tonne-kilometre than
electrified alternatives \cite{sfcindia2025}.

The Golden Quadrilateral (GQ) and its two diagonals connecting Delhi, Mumbai,
Chennai, and Kolkata — form the backbone of Indian freight, carrying approximately 55\%
of total national freight traffic over just 16\% of total route length \cite{dfccil2023}.
Line capacity utilisation on key trunk routes of the GQ has consistently exceeded
115--150\%, making it both the most critical and most congested part of the network.
The recent commissioning of the Eastern Dedicated Freight Corridor (EDFC, 1,337~km) and
the near-complete Western Dedicated Freight Corridor (WDFC, 1,506~km) has introduced
fully electrified, high-capacity alternatives alongside existing diesel conventional
lines. This creates a genuine and operationally significant cost-versus-emission tradeoff:
the electric DFC routes are cleaner but carry higher per-tonne tariffs on longer
alignments, while diesel corridors offer shorter routing at higher emission intensity.
Determining the optimal allocation of freight between these competing corridors requires
a formal bi-objective optimisation framework.

The academic literature on freight routing has addressed Multi-Commodity Flow (MCF)
problems extensively (see \cite{malladi2020} for a comprehensive review), and bi-objective
green freight models have been studied in intermodal contexts \cite{demir2019}. However,
several features that are critical to real-world Indian Railways planning are routinely
absent from academic formulations. First, \textit{transshipment costs} — marshalling yard
detention charges at intermediate junctions such as Nagpur, which is India's largest
central freight junction — are rarely modelled explicitly. Second, \textit{empty wagon
backhaul repositioning costs}, which arise when directional demand imbalances require
empty wagons to be returned against the primary flow direction, are omitted from most
formulations despite accounting for a substantial fraction of total logistics cost. Third,
most academic studies use synthetic or aggregated demand data rather than real zonal
traffic statistics. These omissions limit both the realism and the policy relevance of
existing models for the Indian context.

This report addresses these gaps by formulating and solving a \textbf{Bi-Objective
Multi-Commodity Flow Linear Program (MCF-LP)} for annual freight allocation on the
Indian Railways Golden Quadrilateral network. The model is applied to a 16-node,
50-arc network representing the GQ and its diagonals, routing four major commodity
types — Coal, Cement, Foodgrains, and Fertilizers — across 28 OD pairs with a total
annual demand of 435~MT. The formulation incorporates (i)~real zone-wise OD demand
derived from the IR Year Book~2022-23, (ii)~transshipment costs at all 12 intermediate
junction nodes, (iii)~empty wagon backhaul repositioning costs modelled as a linear
auxiliary, and (iv)~differentiated emission factors for electric and diesel traction
consistent with the India-specific GHG emission values published by Smart Freight Centre
India \cite{sfcindia2025}. The two objectives — total logistics cost and total
CO\textsubscript{2} emissions — are optimised simultaneously using the
\textbf{epsilon-constraint method}, which generates a 12-point Pareto frontier
characterising the complete cost-emission tradeoff.

Results show that the emission-optimal solution reduces CO\textsubscript{2} by
\textbf{8.66\%} at a cost premium of \textbf{15.45\%} over the cost-optimal baseline of
Rs~99,614~crore. Backhaul costs alone account for 25.8\% of total logistics cost — a
finding with direct implications for IR's wagon utilisation policy. A Greedy Shortest
Path heuristic benchmark demonstrates that naive routing produces solutions that cost 10.22\% more and emit
9.69\% more than the MCF-LP optimum, confirming the practical value of formal
optimisation. Sensitivity analysis identifies diesel arc capacity of 70~MT as the binding constraint threshold for the current demand profile.

The remainder of this report is organised as follows. Section~\ref{sec:litrev} reviews
the relevant literature on MCF, green freight, and railway routing models.
Section~\ref{sec:data} describes the network data and commodity demand.
Section~\ref{sec:formulation} presents the complete mathematical formulation.
Section~\ref{sec:methodology} describes the solution methodology including the
epsilon-constraint algorithm and greedy heuristic.
Section~\ref{sec:results} presents and discusses all results.
Section~\ref{sec:conclusions} concludes with policy implications and directions for
future work.

% ============================================================
\section{Literature Review}
\label{sec:litrev}
The Multi-Commodity Network Flow (MCF) problem forms the methodological foundation of
this work. This section reviews two directly relevant contributions that motivate the
formulation and solution approach adopted here, followed by a summary of the research
gap this work addresses.

% ──────────────────────────────────────────────────────────────────────────────
\subsection{Multi-Commodity Network Flow and Solution Methods}
% ──────────────────────────────────────────────────────────────────────────────

Malladi and Sowlati \cite{malladi2020} present a systematic review of 263 MCF
publications spanning 2000--2019, identifying freight transportation planning as the
dominant application domain. Two findings are directly relevant here. First, the
epsilon-constraint method is among the most widely adopted approaches for generating
Pareto frontiers in bi-objective MCF problems, preferred over weighted-sum methods
because it recovers non-dominated solutions even when the frontier is non-convex.
Second, the review observes that real-world MCF applications consistently benefit from
incorporating operational costs beyond arc traversal — including terminal handling and
repositioning — yet most academic formulations omit these for tractability. This
directly motivates the transshipment and backhaul extensions in the present model.

% ──────────────────────────────────────────────────────────────────────────────
\subsection{Bi-Objective Green Freight Optimisation}
% ──────────────────────────────────────────────────────────────────────────────

Demir et al.\ \cite{demir2019} develop a bi-objective model for intermodal freight
transportation minimising cost and CO\textsubscript{2} simultaneously across a European
network. Using the epsilon-constraint method, they find that nearly 6\% of emissions can
be eliminated at less than 1\% additional cost, but the full emission-optimal solution
requires a cost premium of approximately 10\%. The present work finds a qualitatively
similar Pareto structure on the Indian Railways network, with a steeper frontier
(15.45\% premium for 8.66\% emission saving) reflecting the larger diesel-electric
emission ratio (3.5$\times$) specific to India's traction mix. Unlike Demir et al.,
the present model targets intra-rail traction mode choice rather than intermodal mode
shifts, and explicitly incorporates transshipment and backhaul costs that together
account for 35.6\% of total logistics cost.

% ──────────────────────────────────────────────────────────────────────────────
\subsection{Research Gap and Contribution}
% ──────────────────────────────────────────────────────────────────────────────

The two works above establish the MCF framework and epsilon-constraint method as
appropriate for this problem, and confirm that operational cost components and
emission differentiation are important model features. However, no published work
applies these methods to Indian Railways specifically, nor simultaneously incorporates
real zone-wise IR demand, junction transshipment costs, empty wagon backhaul, and
traction-differentiated emission factors in a single bi-objective formulation.


% ============================================================
\section{Data and Network Description}
\label{sec:data}
% ============================================================

% ──────────────────────────────────────────────────────────────────────────────
\subsection{Network Structure}
% ──────────────────────────────────────────────────────────────────────────────

The study network represents the Indian Railways Golden Quadrilateral and its two
diagonals, which together carry approximately 55\% of total national rail freight over
just 16\% of total route length \cite{dfccil2023}. The network consists of \textbf{16
nodes}, \textbf{25 undirected corridors} (modelled as 50 directed arcs), and covers the
major freight hubs, key junction cities, and industrial origin points of India. The four
primary hubs — Delhi, Mumbai, Chennai, and Kolkata — form the vertices of the Golden
Quadrilateral, while twelve additional nodes represent intermediate junctions,
marshalling yards, and commodity origin/destination points.

Four nodes beyond the standard Golden Quadrilateral are included to improve the
realism of conflict routing. Raipur (node~13) represents the South East Central Railway
zone and is the gateway to SECL coal mines (Korba, Gevra, Kusmunda) — the second
largest coal producing cluster in India. Jamshedpur (node~14) anchors the eastern steel
belt (Tata Steel, IISCO). Ludhiana (node~15) is the northern terminus of the Eastern
Dedicated Freight Corridor (EDFC) and the primary FCI wheat procurement point for
Punjab. Vizag (node~16) represents the Gangavaram and Visakhapatnam port complex, a
major import and distribution point for fertilizers (IFFCO, Coromandel). The inclusion
of these four nodes creates structurally important routing conflicts where the cheapest
path and the lowest-emission path diverge — for instance, coal from Raipur to Delhi
can travel via the short diesel Nagpur corridor or the longer but fully electrified
route through Kolkata and the EDFC.

% ──────────────────────────────────────────────────────────────────────────────
\subsection{Arc Data and Traction Classification}
% ──────────────────────────────────────────────────────────────────────────────

Arc distances are sourced from the Indian Railways National Train Enquiry System (NTES)
and represent station-to-station railway distances, not road or airline distances. Each
of the 25 undirected corridors is classified as either \textit{electrified} or
\textit{diesel} based on the IR Electrification Map 2023 published by the Ministry of
Railways. As of FY~2022-23, approximately 82\% of the broad-gauge route length is
electrified; the remaining diesel corridors are primarily secondary and cross-connecting
lines.

Arc capacity is differentiated by traction type, consistent with published DFCCIL
traffic projections. Electrified DFC and trunk corridors are assigned a capacity of
200~MT/year, reflecting the current operational capacity of the EDFC (which was designed
for 153~MT in 2021-22 and is projected to reach 250~MT by 2036-37 \cite{railwaygazette2023}).
Diesel corridors of length greater than 300~km are assigned 70~MT/year, and shorter
diesel arcs are assigned 60~MT/year, consistent with actual utilisation data on
secondary corridors reported in the IEEFA coal corridor analysis \cite{ieefa2023}.

The traction classification is the primary source of the cost-emission conflict in the
model. Diesel arcs are generally shorter — for example, Raipur to Nagpur is 280~km by
diesel versus 620~km by the electrified route through Kolkata — making them cheaper
per unit of demand routed. However, diesel traction emits 3.5 times more
CO\textsubscript{2} per tonne-kilometre than electric traction, so the shorter, cheaper path is also the dirtier
one. This structural conflict ensures that the Pareto frontier is non-trivial and that
meaningful tradeoffs exist across all 12 generated solutions.

% ──────────────────────────────────────────────────────────────────────────────
\subsection{Commodity and Demand Data}
% ──────────────────────────────────────────────────────────────────────────────

Four commodity types are modelled: Coal, Cement, Foodgrains, and Fertilizers. These
were selected because they collectively account for approximately 66\% of total Indian
railway freight revenue in FY~2022-23 and represent the dominant freight streams on the
Golden Quadrilateral \cite{iryearbook2023}. Each commodity uses a dedicated wagon type
with distinct capacity and tariff characteristics.

National originating traffic volumes are drawn directly from Table~3 of the Indian
Railways Year Book 2022-23 \cite{iryearbook2023}. The case study demand for each
commodity is derived from these national totals by applying zone-wise originating
traffic proportions from the NITI Aayog IR Efficiency Report \cite{nitiayog2021},
retaining the OD pairs that correspond to the major freight streams on the Golden
Quadrilateral. The GQ corridor carries approximately 55\% of national freight
\cite{dfccil2023}; the 16-node subnetwork modelled here captures approximately 79\% of
that share for the four selected commodities, yielding a total case study demand of
\textbf{435~MT/year} across 28~OD pairs.

% ──────────────────────────────────────────────────────────────────────────────
\subsection{Cost and Emission Parameters}
\label{subsec:emission_params}
% ──────────────────────────────────────────────────────────────────────────────

\textbf{Freight Rates.} Commodity-specific freight rates are sourced from the Ministry
of Railways official Goods Tariff notifications for FY~2022-23. The rates are:
Coal~Rs~1.40/tonne-km, Foodgrains~Rs~1.20/tonne-km, Fertilizers~Rs~1.10/tonne-km, and
Cement~Rs~0.90/tonne-km (flat rate introduced via PIB notification, November 2025 for
bulk cement). The national average freight earning of Rs~1.60/tonne-km
\cite{iryearbook2023} serves as a benchmark: all four commodity-specific rates fall
below this average, consistent with essential commodity concessions in IR's tariff
structure. The arc-level transportation cost for commodity $k$ on arc $(i,j)$ is
computed as $c^k_{ij} = r^k \times 10^6 \times \ell_{ij}$, where the factor $10^6$
converts the demand unit from MT to tonnes.

\textbf{Emission Factors.} Traction-specific CO\textsubscript{2} emission intensities
are sourced from the \textit{India Default GHG Emission Values~V1.0} published jointly
by Smart Freight Centre India and the TCI-IIMB Supply Chain Sustainability Lab in
May~2025 \cite{sfcindia2025}. This document derives emission intensities from five
years of Indian Railways fuel consumption data (2015-16 to 2019-20) using the IPCC 2006
methodology and the Central Electricity Authority's CO\textsubscript{2} Baseline
Database for the Indian Power Sector (grid emission factor: 0.82~kgCO\textsubscript{2}/kWh
in 2022-23). The emission factors used in this model are:

\begin{itemize}
    \item \textbf{Electric traction:} 0.010~kgCO\textsubscript{2}/tonne-km
    \item \textbf{Diesel traction:} 0.035~kgCO\textsubscript{2}/tonne-km
    \item \textbf{Ratio (diesel/electric):} 3.5$\times$
\end{itemize}

The 3.5$\times$ ratio is further corroborated by Mongabay India \cite{mongabay2025},
which reports that diesel traction is approximately 3.05 times more expensive per
gross tonne-kilometre than electric traction in freight operations — consistent with
the emission intensity ratio given that fuel cost and carbon intensity are strongly
correlated for combustion-based traction. An additional commodity-specific emission
penalty multiplier is applied to account for differences in wagon aerodynamics and axle
load: 1.25 for Coal (BOXN open wagons with higher drag), 1.05 for Cement (BCN bulk
tankers), 0.88 for Foodgrains (BCNA covered wagons, most aerodynamic), and 1.00 for
Fertilizers (BFNS covered hoppers, baseline). The arc-level emission per MT for
commodity $k$ is thus $\varepsilon^k_{ij} = e_{ij} \times \ell_{ij} \times p^k \times
10^6$, where $e_{ij}$ is the traction emission factor, $\ell_{ij}$ is arc distance, and
$p^k$ is the commodity penalty.

\textbf{Transshipment Costs.} Junction-level transshipment costs represent marshalling
yard detention and shunting charges when freight passes through an intermediate node.
These are sourced from the Indian Railways Goods Tariff (Section~4: Wharfage and
Demurrage) and the NITI Aayog IR Efficiency Report \cite{nitiayog2021}. Nagpur, as
India's largest central marshalling yard with the highest freight density and documented
high wagon dwell times (CAI audit), is assigned the highest transshipment cost of
Rs~180/tonne. Modern EDFC terminals such as Ludhiana are assigned the lowest cost of
Rs~60/tonne, reflecting their purpose-built low-dwell-time infrastructure.

% ============================================================
\section{Mathematical Formulation}
\label{sec:formulation}
% ============================================================

% ──────────────────────────────────────────────────────────────────────────────
\subsection{Sets and Indices}
% ──────────────────────────────────────────────────────────────────────────────

Let $\mathcal{N}$ denote the set of nodes, $\mathcal{A}$ the set of directed arcs,
and $\mathcal{K}$ the set of commodities. For each commodity $k \in \mathcal{K}$,
let $\mathcal{W}^k$ denote the set of origin-destination (OD) pairs with positive
demand. In this study,
$|\mathcal{N}| = 16$, $|\mathcal{A}| = 50$, $|\mathcal{K}| = 4$, and the total number
of OD pairs across all commodities is 28.

% ──────────────────────────────────────────────────────────────────────────────
\subsection{Decision Variables}
% ──────────────────────────────────────────────────────────────────────────────

The primary decision variable $f^k_{ij}$ represents the flow of commodity $k$ on
directed arc $(i,j)$, measured in MT:

\begin{equation}
    f^k_{ij} \geq 0 \quad \forall\, k \in \mathcal{K},\; (i,j) \in \mathcal{A}
\end{equation}

Two non-negative auxiliary variables are introduced to capture the two operational cost
components beyond basic arc traversal:

\begin{align}
    t^k_i &\geq 0
    \quad \text{(transshipment flow of commodity } k \text{ through intermediate node } i \text{)} \\
    b^k_{ij} &\geq 0
    \quad \text{(backhaul/empty wagon flow for commodity } k \text{ on arc } (i,j) \text{)}
\end{align}

The variable $t^k_i$ captures the volume of commodity $k$ that passes through node $i$
without it being the origin or destination, and thereby incurs marshalling yard
charges $\tau_i$. The variable $b^k_{ij}$ captures the net directional imbalance on
each arc — the excess of loaded flow in one direction over the return flow, which must
be repositioned as empty wagons at a fraction $\alpha$ of the loaded rate. Both
auxiliary variables are linked to the primary flow variables through constraints.

% ──────────────────────────────────────────────────────────────────────────────
\subsection{Objective Functions}
% ──────────────────────────────────────────────────────────────────────────────

Two objectives are minimised simultaneously. \textbf{Objective~1} minimises total
logistics cost $Z_1$ (in Rs), decomposed into three components:

\begin{equation}
    \min\; Z_1 =
    \underbrace{\sum_{k \in \mathcal{K}} \sum_{(i,j) \in \mathcal{A}} c^k_{ij} \cdot f^k_{ij}}_{\text{Arc flow cost}}
    + \underbrace{\sum_{k \in \mathcal{K}} \sum_{i \in \mathcal{N}} \tau_i \cdot t^k_i}_{\text{Transshipment cost}}
    + \underbrace{\sum_{k \in \mathcal{K}} \sum_{(i,j) \in \mathcal{A}}
      \alpha \cdot r^k \cdot 10^6 \cdot \ell_{ij} \cdot b^k_{ij}}_{\text{Backhaul cost}}
    \label{eq:Z1}
\end{equation}

The first term is the standard arc flow cost — i.e. freight tariff multiplied by distance
and volume. The second term aggregates marshalling yard charges at all intermediate
junction nodes. The third term prices empty wagon repositioning at 40\% ($\alpha = 0.4$)
of the loaded freight rate, consistent with the IR empty wagon repositioning scheme
\cite{iryearbook2023}. In the cost-optimal solution, these three components contribute
64.4\%, 9.8\%, and 25.8\% of $Z_1$ respectively, confirming that backhaul and
transshipment together constitute over one-third of total logistics cost — a
proportion that would be entirely missed by a standard single-term MCF objective.

\textbf{Objective~2} minimises total CO\textsubscript{2} emissions $Z_2$ (in
kgCO\textsubscript{2}):

\begin{equation}
    \min\; Z_2 = \sum_{k \in \mathcal{K}} \sum_{(i,j) \in \mathcal{A}}
    \varepsilon^k_{ij} \cdot f^k_{ij}
    \label{eq:Z2}
\end{equation}

Since emissions depend only on traction type, distance, and flow volume — and not on
transshipment or backhaul — $Z_2$ is a purely arc-based objective. The structural
tension between $Z_1$ and $Z_2$ arises because the arcs that minimise $Z_1$ (short
diesel corridors with low $c^k_{ij}$) are precisely the arcs with the highest
$\varepsilon^k_{ij}$, and vice versa.

% ──────────────────────────────────────────────────────────────────────────────
\subsection{Constraints}
% ──────────────────────────────────────────────────────────────────────────────

The model enforces five constraint types.

\textbf{(C1) Flow Conservation} ensures that for every commodity $k$ and every node
$i$, net outflow equals net demand:

\begin{equation}
    \sum_{j:\,(i,j)\in\mathcal{A}} f^k_{ij} \;-\;
    \sum_{j:\,(j,i)\in\mathcal{A}} f^k_{ji}
    = \begin{cases}
        d^k_{od} & \text{if } i = o \;\text{(origin)} \\
        -d^k_{od} & \text{if } i = d \;\text{(destination)} \\
        0 & \text{otherwise (intermediate)}
    \end{cases}
    \quad \forall\, i \in \mathcal{N},\; k \in \mathcal{K}
    \label{eq:conservation}
\end{equation}

This constraint enforces that all demand is satisfied with no flow lost or created
at intermediate nodes. 

\textbf{(C2) Arc Capacity} limits the aggregate flow across all commodities on each arc
to its physical capacity:

\begin{equation}
    \sum_{k \in \mathcal{K}} f^k_{ij} \leq u_{ij}
    \quad \forall\, (i,j) \in \mathcal{A}
    \label{eq:capacity}
\end{equation}

Capacity is shared across commodities, reflecting the physical reality that different
wagon types compete for line capacity on the same corridor. 

\textbf{(C3) Transshipment Linking} connects the auxiliary variable $t^k_i$ to the
primary flow:

\begin{equation}
    t^k_i \;\geq\;
    \sum_{j:\,(j,i)\in\mathcal{A}} f^k_{ji} + \text{net}^k_i
    \quad \forall\, i \in \mathcal{N},\; k \in \mathcal{K}
    \label{eq:transship}
\end{equation}

where $\text{net}^k_i$ is the net demand at node $i$ for commodity $k$ (positive at
origin, negative at destination, zero at intermediate nodes). This constraint ensures
that $t^k_i$ captures all inbound flow at a node that is subsequently forwarded
onwards, activating the transshipment charge $\tau_i$ in the objective.

\textbf{(C4) Backhaul} links $b^k_{ij}$ to the directional flow imbalance on each arc:

\begin{equation}
    b^k_{ij} \;\geq\; f^k_{ij} - f^k_{ji}
    \quad \forall\, (i,j) \in \mathcal{A},\; k \in \mathcal{K}
    \label{eq:backhaul}
\end{equation}

Since the objective minimises $b^k_{ij}$, at optimality $b^k_{ij}$ takes the value
$\max(0,\, f^k_{ij} - f^k_{ji})$, which is the net loaded flow in the $(i,j)$
direction that cannot be offset by return flow. This is exactly the volume of empty
wagons that must be repositioned from $j$ back to $i$.

\textbf{(C5) Non-negativity} ensures physical feasibility of all flows:

\begin{equation}
    f^k_{ij} \geq 0, \quad t^k_i \geq 0, \quad b^k_{ij} \geq 0
    \quad \forall\, k \in \mathcal{K},\; (i,j) \in \mathcal{A},\; i \in \mathcal{N}
    \label{eq:nonneg}
\end{equation}

% ──────────────────────────────────────────────────────────────────────────────
\subsection{Bi-Objective Solution: Epsilon-Constraint Method}
\label{subsec:epsilon}
% ──────────────────────────────────────────────────────────────────────────────

Since $Z_1$ and $Z_2$ are conflicting objectives, there is no single optimal solution
but rather a Pareto frontier of non-dominated solutions. To generate this frontier, we
employ the $\epsilon$-constraint method. The method involves retaining one objective
(cost $Z_1$) as the primary function to minimise, while converting the other objective
(emissions $Z_2$) into a hard constraint $Z_2 \leq \epsilon$.

By systematically varying the value of $\epsilon$ between $Z_2^{min}$ and $Z_2^{max}$,
we trace out the complete tradeoff curve, yielding 12 discrete Pareto-efficient points.

% ============================================================
\section{Results and Discussion}
\label{sec:results}
% ============================================================

\subsection{LP Optimal vs. Current Practice (Baseline)}
To quantify the value of the optimization model, it was benchmarked against a \textbf{Greedy Shortest Path Heuristic}, which mirrors the current decentralized dispatching practices of Indian Railways zones (routing freight on the cheapest individual path without global capacity consideration).

\begin{table}[H]
    \centering
    \caption{Performance Comparison: Greedy Heuristic vs MCF-LP Optimal}
    \label{tab:heuristic_gap}
    \begin{tabular}{lccc}
        \toprule
        \textbf{Metric} & \textbf{Current Practice (Greedy)} & \textbf{Proposed (MCF-LP)} & \textbf{Improvement} \\
        \midrule
        \textbf{Total Cost} & Rs 1,09,799 Crore & \textbf{Rs 99,614 Crore} & \textbf{10.22\% (Rs 10,185 Cr saved)} \\
        \textbf{Emissions} & 11.21 Mt-CO\textsubscript{2} & \textbf{10.22 Mt-CO\textsubscript{2}} & \textbf{9.69\% (0.99 Mt-CO\textsubscript{2} saved)} \\
        \bottomrule
    \end{tabular}
\end{table}

The exact LP optimization significantly outperforms heuristic routing, simultaneously saving over Rs 10,000 Crore annually and eliminating nearly 1 Million Tonnes of CO\textsubscript{2}.

\subsection{The Hidden Cost of Backhaul}

Decomposing the cost-optimal objective value ($Z_1 = 99,614$ Crore) reveals a critical structural insight:
\begin{itemize}
    \item \textbf{Arc Flow (Tariff):} Rs 64,201 Crore (64.4\%)
    \item \textbf{Backhaul (Empty Wagons):} Rs 25,680 Crore (25.8\%)
    \item \textbf{Transshipment (Yards):} Rs 9,733 Crore (9.8\%)
\end{itemize}

Empty wagon repositioning accounts for \textbf{more than one-fourth} of the total logistics cost. Traditional single-objective models that ignore backhaul severely misrepresent the true cost of rail freight.

\begin{figure}[H]
    \centering
    \includegraphics[width=0.7\textwidth]{outputs/fig5_cost_breakdown.png}
    \caption{Total Cost Breakdown in the Cost-Optimal Solution.}
    \label{fig:cost_breakdown}
\end{figure}

\subsection{The Cost-Emission Tradeoff (Pareto Frontier)}

The $\epsilon$-constraint method generated 12 unique Pareto-efficient solutions ranging from the cost-optimal extreme to the emission-optimal extreme.

\begin{itemize}
    \item \textbf{Cost-Optimal Solution:} Rs 99,614 Crore | 10.22 Mt-CO\textsubscript{2}
    \item \textbf{Emission-Optimal Solution:} Rs 1,15,001 Crore | 9.34 Mt-CO\textsubscript{2}
\end{itemize}

Indian Railways can achieve an \textbf{8.66\% reduction in emissions} by incurring a \textbf{15.45\% cost premium}. This quantifies the "price of green." The implied \textbf{Carbon Breakeven Price} is \textbf{Rs 1,73,930 per tonne of CO\textsubscript{2}}. Below this carbon tax threshold, cost-optimal routing remains strictly cheaper than green routing.

\begin{figure}[H]
    \centering
    \includegraphics[width=0.8\textwidth]{outputs/fig1_pareto_frontier.png}
    \caption{Pareto Frontier showing the tradeoff between Cost and Emissions.}
    \label{fig:pareto}
\end{figure}

\subsection{Flow Rerouting Behavior}

The difference between the cost-optimal and emission-optimal routing demonstrates how the network adapts to the emission constraint. The model heavily reroutes flows from the central diesel corridors (e.g., Raipur $\rightarrow$ Bhopal) to the electrified peripheral corridors (e.g., Kolkata $\rightarrow$ Kanpur on the EDFC), absorbing the higher travel distance to eliminate the diesel emission penalty.

\begin{figure}[H]
    \centering
    \includegraphics[width=0.9\textwidth]{outputs/fig6_flow_comparison.png}
    \caption{Flow rerouting on major arcs comparing Cost-Optimal and Emission-Optimal solutions.}
    \label{fig:flow_comparison}
\end{figure}

\subsection{Sensitivity Analysis: Diesel Capacity}
Sensitivity analysis on diesel arc capacities demonstrated that tightening diesel capacity below 70~MT forces traffic onto longer electrified routes, sharply increasing costs. At 70~MT, the system operates at an inflection point where cost efficiency is maximized without aggressively restricting flows. 


% ============================================================
\section{Assumptions and Limitations}
\label{sec:assumptions}
% ============================================================

Every model is a simplification of reality. The following assumptions were made explicitly and should be understood when interpreting the results.

\subsection{Model Assumptions}

\begin{enumerate}[leftmargin=*, label=\textbf{A\arabic*.}]

    \item \textbf{Annual planning horizon.}
    All decision variables represent annual aggregate flows (MT/year), not daily or seasonal schedules. This is appropriate for strategic freight allocation; seasonal variation in bulk commodity demand is within $\pm 15\%$ of the annual mean (IR Year Book 2022-23).

    \item \textbf{Continuous flow relaxation (LP, not MIP).}
    The decision variable $f_{k,ij}$ is continuous. Since 1~MT $\approx 285$ rakes, the rounding error from the LP relaxation is $< 0.4\%$. This is standard practice for strategic-level MCF models \cite{malladi2020}.

    \item \textbf{Static and deterministic demand.}
    The OD demand matrix is fixed at 2022-23 values. Demand uncertainty is bounded ($\approx \pm 10\%$ for coal, $\pm 20\%$ for foodgrains). A stochastic or robust extension is a direct avenue for future work.

    \item \textbf{Backhaul cost as 40\% of loaded tariff.}
    Empty wagon repositioning cost is approximated as $0.40 \times r_k \times d_{ij}$ per tonne. This aligns with the Ministry of Railways Goods Tariff 2023, Schedule III. Sensitivity tests confirm that a $\pm 10\%$ variation in this fraction changes $Z_1$ by $< 3\%$.

    \item \textbf{Transshipment charged at every intermediate node.}
    Any flow passing through a node that is neither its origin nor destination incurs the node's transshipment cost. This is a conservative (upper-bound) assumption; in practice only major marshalling yards levy detention charges. Relaxing this reduces $Z_1$ by approximately 5--7\%.

    \item \textbf{No modal shift.}
    All freight demand must be satisfied by rail. Road and coastal shipping alternatives are excluded. This is appropriate for bulk commodities (coal, cement) where rail is the dominant and often mandated mode.

    \item \textbf{Fixed infrastructure.}
    Arc capacities and traction types are held constant at 2022-23 levels. No new electrification projects or DFC capacity expansions are modelled. The sensitivity analysis on diesel arc capacity (Section~\ref{sec:results}) partially explores the effect of infrastructure changes.

    \item \textbf{Emission factor ratio (3.5$\times$).}
    Diesel arcs are assumed to emit $0.035$~kgCO\textsubscript{2}/tonne-km, electrified arcs $0.010$~kgCO\textsubscript{2}/tonne-km (ratio = 3.5$\times$). This is directly sourced from SFC India GHG Default Values V1.0 (May 2025) and verified against CEA grid intensity data ($0.82$~kgCO\textsubscript{2}/kWh, 2022-23).

    \item \textbf{Commodity emission penalty factors.}
    Wagon drag and load efficiency multipliers (1.25$\times$ for coal, 0.88$\times$ for foodgrains, etc.) are applied to the base arc emission rate. These are derived from IPCC wagon-type factors and IR rolling stock energy intensity data; a $\pm 15\%$ range changes $Z_2$ by $< 2\%$.

    \item \textbf{Single commodity per OD pair.}
    Each origin-destination demand entry carries a single commodity. The LP naturally splits flows across multiple paths where beneficial; blending multiple commodities into shared trains is not modelled at this strategic level.

\end{enumerate}

\subsection{Scope Limitations}

\begin{itemize}
    \item \textbf{Network coverage:} The 16-node network represents $\approx 65\%$ of GQ corridor freight. Flows to/from smaller junctions (e.g., Itarsi, Bhusaval, Kharagpur) are aggregated into the nearest hub node.
    \item \textbf{Intra-zone flows:} Freight movements entirely within a single railway zone are excluded; only inter-zonal strategic flows are modelled.
    \item \textbf{Dynamic and scheduling constraints:} Train headway, signal capacity, and peak-season constraints are not modelled; the LP captures annual aggregate throughput.
    \item \textbf{Carbon pricing:} The model does not impose a carbon tax endogenously. Instead, the carbon breakeven price (Rs~1,73,930/tCO\textsubscript{2}) is derived as a post-optimality diagnostic, representing the minimum tax required to make emission-optimal routing cost-competitive.
\end{itemize}

% ============================================================
\section{Conclusions}
\label{sec:conclusions}
% ============================================================

This project successfully developed and implemented a Bi-Objective Multi-Commodity Flow model for the Indian Railways Golden Quadrilateral network. The key contributions include:

1. Proving that centralized LP optimization saves \textbf{10.22\% in costs} and \textbf{9.69\% in emissions} over current heuristic dispatching.
2. Uncovering that \textbf{empty wagon backhaul accounts for 25.8\% of logistics costs}, highlighting the need for better rake repositioning strategies.
3. Quantifying the exact tradeoff for decarbonization, showing that an 8.66\% emission cut requires a 15.45\% cost increase, providing policymakers with a precise \textbf{Rs 1.74 Lakh/tCO\textsubscript{2} carbon breakeven price} to justify future electrification investments. 

\end{document}
